%% THIS IS THE MAIN FILE %%
% Here the main arrangement of your thesis is determined. 
% In this file you read-in all chapters (as separate .tex files)
% This is also the file you 'Build' with a LaTeX editor, like TeXmaker.
% The .cls file (MScThesis.cls) contains all details regarding copyrights and so on. Take a look and adjust where applicable. But be careful changing this file: it determines the whole lay-out of your thesis!
% This MSc thesis standard layout is optimized for double sided printing on A4 format paper.
% Good luck and have fun with your Master Research in Applied Geophysics! Kind regards, Niels Grobbe

\documentclass[11pt,a4paper,openany,oneside]{MScThesis}


\usepackage[
  style=apa,
  backend=biber,
  natbib=true
]{biblatex}
\addbibresource{master_thesis.bib} % Points to your .bib database
%
%\usepackage{pifont}
\mscName{Salome Anna Bachmann}
\mscDate{\today}
\mscTitle{Test}
\mscSubTitle{Subtitle}
\mscKeyWords{thesis, msc, subject}
%
%\mscBackPicture{2660PF3}    % eps of 21 * 29.7 cm
\mscReaderOne{Name of first supervisor}
\mscReaderTwo{Name of second supervisor}
\mscReaderThree{Name of third supervisor or committee member}
%\mscReaderFour{}
%
\setThesisInfo

%
\begin{document}
%
%============================= Front matter ========================================
\frontmatter %
%
% Make a hell of a lot of title pages
    \maketitle

%============ ! ! ! IMPORTANT FOR MSC THESES AT RWTH AACHEN ONLY ! ! ! ==============

% In accordance with the comprehensive exam regulations of the RWTH Aachen students writing their master thesis at the RWTH must sign an affidavit for written exam work, confirming that the exam performance was completed without unauthorized assistance and that no other sources or aids were used than the ones listed in the written work.

% For your thesis you must use the official affidavit in German "Formular_Eidesstattliche_Versicherung_neu.pdf" which is included in the MSc_layout folder. This version is legally binding; the English translation "Statutory_Declaration_in_Lieu_of_an_Oath.pdf" is provided for your information only.

% INSTRUCTIONS

%1 Open the form "Formular_Eidesstattliche_Versicherung_neu.pdf" with your pdf reader (Adobe Reader recommended) and fill it in. Please do not delete the blank page.

%2 Uncomment the command below:

% \includepdf[scale=1,pages=-]{RWTH_only_Eidesstattliche_Versicherung.pdf}

%3 Sign the form once your thesis is printed

%====================================================================================

% Abstract
    \nonumchap{Abstract}



    \cleardoublepage
%
% Acknowledgements
    \nonumchap{Acknowledgements}%
   
    \vspace*{15mm}

    \noindent
    ETH Zürich \hfill \mscname\\ % the university where you carried out your final MSc thesis research
    \mscdate

%
% table of contents, (\toc or \toclof or \tocloflot )
   \tocloflot 

%
% Nomenclature
    \printnomencl %
%
% Acronyms
    \nonumchap{Acronyms} %
    %
    \begin{acronym}%
        \acro{DAS} Distributed Acoustic Sensing
        \acro{MPM} Material Point Method
        \acro{DEM} Discrete Element Method
        \acro{FE, FEM} Finite Element Method
        \acro{PST} Propagation Saw Test
    \end{acronym}%
    %
    \newpage%
%
%
%============================= Main matter =========================================
%
\mainmatter
%
% Introduction
\include{chapter1} % notice how the reference to the introduction takes place.. It refers to the name.tex
\nomenclature{$\boldsymbol{t_o}$}{Subfault onset time [s], e.g. $t_o=(x-x_{\rm start})/v_{\rm rupture}$.}
\nomenclature{$\boldsymbol{x},\,\boldsymbol{x}_{\rm start}$}{Horizontal coordinate and rupture starting position [m].}
\nomenclature{$\boldsymbol{v}_{\rm rupture}$}{Prescribed rupture propagation speed [m s$^{-1}$].}
\nomenclature{$\boldsymbol{t}_{rise}$}{Rise time of a subfault/source [s].}
\nomenclature{$\boldsymbol{\Delta x}$}{Subfault width or discretisation length in x [m].}
\nomenclature{$\boldsymbol{x}_{\rm slip}$}{Prescribed slip on a subfault [m].}
\nomenclature{$\boldsymbol{M}_0$}{Scalar seismic moment per subfault [N m].}
\nomenclature{$\boldsymbol{m}_{ij}$}{Moment tensor components [N m].}
\nomenclature{$\boldsymbol{r}_{\rm mode}$}{Mode mix parameter [dimensionless].}
\nomenclature{$\boldsymbol{l}_{cz}$}{Cohesive process (fracture) zone length [m].}
\nomenclature{$\boldsymbol\tau_s$}{Shear strength of interface [Pa].}
\nomenclature{$\boldsymbol{l}_{ch}$}{Irwin characteristic length [m].}
\nomenclature{$\boldsymbol{\lambda}_{\min}$}{Shortest wavelength to resolve [m].}
\nomenclature{$\boldsymbol{v}_{\min},\,\boldsymbol{v}_{\max}$}{Minimum and maximum relevant wave speeds in the model [m s$^{-1}$].}
\nomenclature{$\boldsymbol{\Delta t}$}{Time step [s] for explicit time integration.}
\nomenclature{$\boldsymbol{\xi}$}{Normalized opening [dimensionless].}
\nomenclature{$\boldsymbol{\delta}$}{Local opening (separation) displacement across the interface [m].}
\nomenclature{$\boldsymbol{\delta}_c$}{Mixed-mode critical opening [m] used to nondimensionalize openings.}
\nomenclature{$\boldsymbol{\delta}_{n,c},\,\boldsymbol{\delta}_{s,c}$}{Single-mode critical openings for Mode I (normal) and Mode II (shear) [m].}
\nomenclature{$\boldsymbol{T}(\boldsymbol{\xi})$}{Cohesive traction as a function of normalised opening [Pa or N m$^{-2}$].}
\nomenclature{$\boldsymbol{T}_{\mathrm{peak}}$}{Peak cohesive traction [Pa].}
\nomenclature{$\boldsymbol{\phi}(\boldsymbol{\Delta})$}{Mixed-mode cohesive potential [J m$^{-2}$] with $\Delta=(\Delta_n,\Delta_t)$.}
\nomenclature{$\boldsymbol{\Delta}_n,\,\boldsymbol{\Delta}_t$}{Normal and tangential separations [m].}
\nomenclature{$\boldsymbol{\delta}_n,\,\boldsymbol{\delta}_t$}{Normal and tangential length scales in the cohesive potential [m].}
\nomenclature{$\boldsymbol{q}$}{Coupling parameter in the mixed-mode cohesive potential [dimensionless].}
\nomenclature{$\boldsymbol{\phi}_n,\,\boldsymbol{\phi}_t$}{Work (energies) of separation for pure normal and tangential modes [J m$^{-2}$].}
\nomenclature{$\boldsymbol{\sigma}_{\max},\,\boldsymbol{\tau}_{\max}$}{Peak normal strength and peak shear strength [Pa].}
\nomenclature{$\boldsymbol{\sigma}_n,\,\boldsymbol{\tau}_s$}{Normal and shear strength [Pa].}
\nomenclature{$\boldsymbol{G}_I,\,\boldsymbol{G}_{II}$}{Fracture energies for Mode I and Mode II [J m$^{-2}$].}
\nomenclature{$\boldsymbol{v}_s,\,\boldsymbol{v}_p$}{Pressure and Shear seismic wave speeds [m s$^{-1}$].}
\nomenclature{$\boldsymbol{\rho}$}{Mass density [kg m$^{-3}$].}
\nomenclature{$\boldsymbol{\mu},\,\boldsymbol{\lambda}$}{Lamé parameters [Pa].}
\nomenclature{$\boldsymbol{E}$}{Young's modulus [Pa].}
\nomenclature{$\boldsymbol{l}_{crack},\,\boldsymbol{x}_0$}{Local crack process length and initial crack front position [m] used in local coordinate definitions.}
\nomenclature{$\boldsymbol{\nabla}$}{Component-wise partial derivative operator (nabla).}
\nomenclature{$\boldsymbol{\theta}$}{Reference coordinate system in respect to mesh elements used in SALVUS.}
\nomenclature{$\boldsymbol{w}$}{Test function used to obtain the weak form of FE equations [dimensionless].}
\nomenclature{$\boldsymbol{G}$}{Model domain [dimensionless].}
\nomenclature{$\boldsymbol{V}$}{Galerkin space [dimensionless].}
\nomenclature{$\boldsymbol{u}$}{Scalar displacement potential [dimensionless].}
\nomenclature{$\boldsymbol{\phi}$}{Set of functions which approximate solution of w or u [dimensionless].}
\nomenclature{$\boldsymbol{F}$}{Forcing term of elastic wave equation.}
\nomenclature{$\boldsymbol{M}$}{Mass term of elastic wave equation.}
\nomenclature{$\boldsymbol{S}$}{Stiffness term of elastic wave equation.}
\nomenclature{$\boldsymbol{B}$}{Boundary term of elastic wave equation.}
\nomenclature{$\boldsymbol{\partial}$}{Partial derivative in respect to one variable.}


\cleardoublepage

%
% First Part
    \part{First Part} % you can divide your thesis into different parts, for example Theory&Modeling in part 1 and real data examples in part 2.
		\include{chapter2}




%
% Bibliography
    % Bibliography
    \cleardoublepage
    \printbibliography % in MyBib.bib you add all your reference information, following the correct format.
    % The *.bin file needs to be compiled using the Bib-compile option of your LaTeX text editor. Sometimes, the bib file needs to be built several times, as well as the main file, before all references occur correctly in your PDF.


%============================= Back matter =========================================
\appendix

    \chapter{The back of the thesis}

    \section{An appendix section}

    \subsection{An appendix subsection with C++ Lisitng}

   

    \subsection{A \matlab $ $ Listing}

 

    \chapter{Yet another appendix}

    \section{Another test section}

    Ok, all is well.

% Index
    \printindex
    \cleardoublepage

\end{document}